\documentclass[12pt]{article}
\usepackage[utf8]{inputenc}
\usepackage[T1]{fontenc}
\usepackage[margin=1in]{geometry}
\usepackage{setspace}
\usepackage{xcolor}
\usepackage[colorlinks=true, urlcolor=blue!60!black, linkcolor=blue!60!black]{hyperref}
\setlength{\parindent}{0em}
\setlength{\parskip}{0.45em}
\renewcommand{\baselinestretch}{1.0}
\begin{document}
\begin{center}
{\Large \textit{New York Times Co. v. United States}, 403 U.S. 713 (1971)\par}
\vspace{0.3em}
{\normalsize Case Brief --- One-Page Edition\par}
\end{center}

\textbf{Memory Jogger.} The Pentagon Papers Case: no prior restraint without meeting a ``heavy burden'' of justification (\href{https://supreme.justia.com/cases/federal/us/403/713/}{Justia}).

\textbf{Facts.} Daniel Ellsberg leaked a classified history of the Vietnam War to the \textit{New York Times} and \textit{Washington Post}; the Nixon Administration sued to stop publication, claiming grave and irreparable injury to national security.

\textbf{Procedural History.} The Second Circuit restrained the \textit{Times}; the D.C. Circuit refused to restrain the \textit{Post}.

\textbf{Votes.} 6--3 \textit{per curiam}, with nine separate opinions.

\textbf{Holding.} Prior restraints carry a heavy presumption of invalidity, and the Government did not meet its burden.

\textbf{Reasoning.} Black and Douglas would bar prior restraints outright; Brennan required proof of inevitable, immediate catastrophe; Stewart required direct, immediate, and irreparable damage; White and Marshall stressed the lack of congressional authorization while leaving prosecution open. The dissenters protested ``unseemly haste'' and would have deferred to executive judgments about foreign affairs.

\textbf{Examples.} \emph{Same-Side:} a leaked historical memo. \emph{Opposite-Side:} a memoir under a review agreement (\textit{Snepp}). \emph{Fence-Sitter:} a statute against naming covert officers.

\textbf{Critique.} Nine opinions left the law unsettled, and prosecution threats may chill more than injunctions.

\textbf{Key Quotation.} ``I am convinced that the Executive is correct with respect to some of the documents involved. But I cannot say that disclosure of any of them will surely result in direct, immediate, and irreparable damage to our Nation or its people. That being so, there can under the First Amendment be but one judicial resolution of the issues before us.'' --- Justice Stewart, 403 U.S. at 730.

\end{document}
